\documentclass[12pt,a4paper]{article}
\usepackage[a4paper,margin=2.5cm]{geometry}
\usepackage{amsmath,amssymb}
\usepackage{icomma} % koma sebagai pemisah desimal, mis. 0,4
\usepackage{enumitem}

\setlength{\parindent}{0pt}
\setlength{\parskip}{6pt}

\begin{document}

\begin{center}
  {\Large\bfseries Tugas 3}\\[2pt]
  {\large Pengantar Probabilitas (STIK4122)}
\end{center}

\begin{tabular}{@{}l@{\;:\;}l}
  Nama  & RIZIQ LILI ULIL ALBAB \\
  NIM   & 058487469 \\
  Prodi & S1 STATISTIKA
\end{tabular}

\section*{Diketahui}

Matriks transisi dengan urutan kondisi Tinggi (T), Sedang (S), dan Rendah (R):
\[
P =
\begin{array}{c@{\;}c}
  & \begin{array}{ccc} \text{T} & \text{S} & \text{R} \end{array} \\
  \begin{array}{c} \text{T} \\ \text{S} \\ \text{R} \end{array} &
  \begin{pmatrix}
    0,4 & 0,3 & 0,3 \\
    0,2 & 0,6 & 0,2 \\
    0   & 0,5 & 0,5
  \end{pmatrix}
\end{array}
\]

Distribusi awal:
\[
\pi_0 = \begin{pmatrix} 0,4 & 0,3 & 0,3 \end{pmatrix}
\]
\textit{Catatan: digunakan $\pi_0 = (0,4 \;\; 0,3 \;\; 0,3)$ sesuai vektor yang tertulis pada soal.}

\section*{1. Distribusi kondisi penjualan pada hari berikutnya}

Rumus: $\pi_1 = \pi_0 P$
\begin{align*}
\pi_1(T) &= (0,4)(0,4) + (0,3)(0,2) + (0,3)(0) = 0,16 + 0,06 + 0 = 0,22 \\
\pi_1(S) &= (0,4)(0,3) + (0,3)(0,6) + (0,3)(0,5) = 0,12 + 0,18 + 0,15 = 0,45 \\
\pi_1(R) &= (0,4)(0,3) + (0,3)(0,2) + (0,3)(0,5) = 0,12 + 0,06 + 0,15 = 0,33
\end{align*}
Jadi,
\[
\boxed{\pi_1 = \begin{pmatrix} 0,22 & 0,45 & 0,33 \end{pmatrix}}
\qquad \text{(cek: } 0,22 + 0,45 + 0,33 = 1\text{)}
\]

\section*{2. Distribusi kondisi penjualan setelah dua hari}

Rumus: $\pi_2 = \pi_1 P = \pi_0 P^2$
\begin{align*}
\pi_2(T) &= (0,22)(0,4) + (0,45)(0,2) + (0,33)(0) = 0,088 + 0,090 + 0 = 0,178 \\
\pi_2(S) &= (0,22)(0,3) + (0,45)(0,6) + (0,33)(0,5) = 0,066 + 0,270 + 0,165 = 0,501 \\
\pi_2(R) &= (0,22)(0,3) + (0,45)(0,2) + (0,33)(0,5) = 0,066 + 0,090 + 0,165 = 0,321
\end{align*}
Sebagai pembanding, matriks transisi dua langkah adalah
\[
P^2 =
\begin{pmatrix}
  0,22 & 0,45 & 0,33 \\
  0,20 & 0,52 & 0,28 \\
  0,10 & 0,55 & 0,35
\end{pmatrix},
\]
dan $\pi_0 P^2$ memberikan hasil yang sama. Jadi,
\[
\boxed{\pi_2 = \begin{pmatrix} 0,178 & 0,501 & 0,321 \end{pmatrix}}
\qquad \text{(cek: } 0,178 + 0,501 + 0,321 = 1\text{)}
\]

\section*{3. Peluang berada pada kondisi Sedang setelah dua hari}

\[
P(X_2 = S) = \pi_2(S) = \boxed{0,501}
\]
Jadi, peluang perusahaan berada pada kondisi penjualan sedang setelah dua hari adalah $0,501$ atau sekitar $50,1\%$.

\section*{4. Distribusi stasioner}

Misalkan $\pi = (\pi_T \;\; \pi_S \;\; \pi_R)$ dengan syarat
\[
\pi = \pi P \qquad \text{dan} \qquad \pi_T + \pi_S + \pi_R = 1.
\]

Dari kolom T:
\begin{align*}
\pi_T &= 0,4\pi_T + 0,2\pi_S + 0\cdot\pi_R \\
0,6\pi_T &= 0,2\pi_S \\
\pi_S &= 3\pi_T \tag{i}
\end{align*}

Dari kolom S:
\begin{align*}
\pi_S &= 0,3\pi_T + 0,6\pi_S + 0,5\pi_R \\
0,4\pi_S &= 0,3\pi_T + 0,5\pi_R \\
0,4(3\pi_T) &= 0,3\pi_T + 0,5\pi_R \\
0,9\pi_T &= 0,5\pi_R \\
\pi_R &= 1,8\pi_T \tag{ii}
\end{align*}

Syarat jumlah sama dengan 1:
\begin{align*}
\pi_T + 3\pi_T + 1,8\pi_T &= 1 \\
5,8\pi_T &= 1 \\
\pi_T &= \frac{1}{5,8} = \frac{5}{29} \approx 0,172
\end{align*}
Sehingga
\[
\pi_S = 3 \cdot \frac{5}{29} = \frac{15}{29} \approx 0,517,
\qquad
\pi_R = 1,8 \cdot \frac{5}{29} = \frac{9}{29} \approx 0,310.
\]

Cek dengan kolom R:
\[
0,3\left(\tfrac{5}{29}\right) + 0,2\left(\tfrac{15}{29}\right) + 0,5\left(\tfrac{9}{29}\right)
= \frac{1,5 + 3 + 4,5}{29} = \frac{9}{29} = \pi_R \quad \checkmark
\]

Jadi, distribusi stasionernya adalah
\[
\boxed{\pi = \begin{pmatrix} \dfrac{5}{29} & \dfrac{15}{29} & \dfrac{9}{29} \end{pmatrix}
\approx \begin{pmatrix} 0,172 & 0,517 & 0,310 \end{pmatrix}}
\]
\end{document}